\documentclass[11pt]{article}
\usepackage{amsmath, amssymb, amsthm}
\usepackage{geometry}
\usepackage{hyperref}
\geometry{letterpaper, margin=1in}

\newtheorem{theorem}{Theorem}[section]
\newtheorem{lemma}[theorem]{Lemma}
\newtheorem{proposition}[theorem]{Proposition}
\theoremstyle{definition}
\newtheorem{definition}{Definition}[section]

\title{Project SOA:\\
Theory of Shared-Basis Attention for BERT-base on Apple M1 Hardware}
\author{Research Dossier}
\date{}

\begin{document}
\maketitle

\begin{abstract}
This paper presents a theoretical framework and methodology for \emph{Project SOA}, which proposes replacing classical Query/Key/Value (Q/K/V) projection matrices in BERT-base attention layers with a \emph{shared global orthogonal basis} constructed via Singular Value Decomposition (SVD). This approach drastically reduces parameter count and computational cost, enabling efficient long-context transformer inference on Apple M1's Unified Memory Architecture and Neural Engine hardware. We provide a comprehensive linear algebraic justification, subspace alignment theory across layers, optimization implications of a fixed shared basis, and layout considerations for rank selection and hardware deployment. This concise paper targets a compact formal theory presentation supplemented by empirical methodology plans.
\end{abstract}

\tableofcontents

\section{Problem Setup and Notation}

Consider BERT-base architecture with:
\begin{itemize}
    \item Number of layers: \(L = 12\)
    \item Hidden dimension: \(d = 768\)
    \item Number of attention heads: \(h = 12\), head size \(d_h = 64\)
    \item Input sequence length \(T\) (targeted for long contexts \(T = 2k \text{ to } 4k\))
\end{itemize}

At each layer \(\ell \in \{1,\dots,L\}\), define the standard attention linear projections:
\[
W_Q^{(\ell)}, W_K^{(\ell)}, W_V^{(\ell)} \in \mathbb{R}^{d \times d}.
\]
Given input hidden states \(X^{(\ell)} \in \mathbb{R}^{T \times d}\),
\[
Q^{(\ell)} = X^{(\ell)} W_Q^{(\ell)}, \quad
K^{(\ell)} = X^{(\ell)} W_K^{(\ell)}, \quad
V^{(\ell)} = X^{(\ell)} W_V^{(\ell)}.
\]

Multi-head attention is fused conceptually into single Q/K/V mappings per layer.

\subsection{Shared-Basis Parameterization}

Project SOA replaces each attention weight matrix by:
\[
W_*^{(\ell)} \approx U A_*^{(\ell)},
\]
where:
\begin{itemize}
    \item \(U \in \mathbb{R}^{d \times r}\) is a global orthonormal basis with rank \(r \ll d\)
    \item \(A_*^{(\ell)} \in \mathbb{R}^{r \times d}\) are small trainable projection matrices
\end{itemize}
Here, \(U\) is fixed (frozen) after estimation, shared across all Q/K/V and all layers, and \(A_*^{(\ell)}\) are adapted in training or derived from pretrained weights.

Two main constructions for \(U\) are considered:
\begin{enumerate}
    \item \textbf{Pretrained-Weight SVD:} SVD on concatenated \(W_Q^{(\ell)}, W_K^{(\ell)}, W_V^{(\ell)}\)
    \item \textbf{Covariance-Based SVD:} SVD on hidden-state covariance matrix across layers/positions
\end{enumerate}

\section{Numerical Linear Algebra Foundations}

\subsection{Eckart–Young–Mirsky Theorem}

\begin{theorem}[Eckart–Young–Mirsky]
Let matrix \(W \in \mathbb{R}^{d \times d}\) have singular value decomposition
\[
W = U_W \Sigma_W V_W^\top,
\]
with singular values \(\sigma_1 \ge \sigma_2 \ge \cdots \ge \sigma_d \ge 0\).

For any \(r < d\), the best rank-\(r\) approximation to \(W\) minimizing the Frobenius norm error is
\[
W_r := U_W^{(:,1:r)} \Sigma_W^{(1:r,1:r)} V_W^{(:,1:r)\top},
\]
and
\[
\min_{\operatorname{rank}(Z) \le r} \|W - Z\|_F = \left( \sum_{i > r} \sigma_i^2 \right)^{1/2} = \|W - W_r\|_F.
\]
\end{theorem}

\begin{proof}[Sketch]
See standard linear algebra references; this follows from properties of SVD truncation preserving optimality in unitarily invariant norms.
\end{proof}

This theorem justifies factorizing attention projections as \(W \approx U A\), where \(U = U_W^{(:,1:r)}\).

\subsection{Shared-Basis Approximation Across Multiple Matrices}

Let \(\{W^{(m)}\}_{m=1}^M\) denote all \(3L\) attention matrices. We seek a single orthonormal basis \(U \in \mathbb{R}^{d \times r}\) minimizing the total approximation error:
\[
\min_{U^\top U = I_r} \sum_{m=1}^M \min_{A^{(m)}} \|W^{(m)} - U A^{(m)}\|_F^2.
\]

\begin{lemma}[Optimal \(A^{(m)}\) for fixed \(U\)]
Given fixed \(U\), the minimizer is \(A^{(m)*} = U^\top W^{(m)}\) with residual error
\[
\| W^{(m)} - U A^{(m)*} \|_F^2 = \| (I - U U^\top) W^{(m)} \|_F^2.
\]
\end{lemma}

\begin{proof}
By orthogonal projection properties and Pythagorean theorem.
\end{proof}

Therefore, the problem reduces to:
\[
\min_{U^\top U = I_r} \sum_{m=1}^M \| (I - U U^\top) W^{(m)} \|_F^2.
\]

Stacking all column vectors of \(W^{(m)}\) into \(\mathcal{X} \in \mathbb{R}^{d \times N}\), where \(N = M \times d\),
this is the classical PCA objective on \(\mathcal{X}\).

\begin{proposition}[Shared-Basis PCA]
The minimizing \(U\) consists of the top \(r\) left singular vectors of \(\mathcal{X}\).
\end{proposition}

\begin{proof}[Sketch]
Equivalent to projection onto best rank \(r\) subspace minimizing reconstruction error.
\end{proof}

\subsection{Global Energy Criterion for Rank Selection}

Define total energy
\[
E_{\text{total}} = \|\mathcal{X}\|_F^2 = \sum_i \sigma_i^2(\mathcal{X}),
\]
and rank-\(r\) preserved energy
\[
E(r) = \sum_{i=1}^r \sigma_i^2(\mathcal{X}).
\]
Choose rank \(r\) s.t.
\[
\eta(r) := \frac{E(r)}{E_{\text{total}}} \ge \tau,
\]
for threshold \(\tau \approx 0.99\).

---

\section{Subspace Alignment and Shared Basis Justification}

\begin{definition}[Subspace Distance]
For two \(r\)-dimensional orthonormal bases \(U^{(1)}, U^{(2)} \in \mathbb{R}^{d \times r}\), define
\[
d_{\mathrm{sub}}(U^{(1)}, U^{(2)}) := \max_{i=1}^r \sin \theta_i,
\]
where \(\{\theta_i\}\) are principal angles between the two subspaces.
\end{definition}

\begin{proposition}[Error bound with shared basis]
If \(W_r^{(m)}\) is the best rank-\(r\) approximation of \(W^{(m)}\) with subspace \(\mathcal{S}^{(m)}\), and \(U\) is a shared basis, then
\[
\|W^{(m)} - U A^{(m)}\|_F^2 \le \| W^{(m)} - W_r^{(m)} \|_F^2 + C_m d_{\mathrm{sub}}^2(\mathcal{S}^{(m)}, \operatorname{span}(U)),
\]
for some constant \(C_m\).
\end{proposition}

\begin{proof}[Sketch]
Employ decomposition of error into components in and orthogonal to shared basis and apply Davis-Kahan-type bounds.
\end{proof}

This motivates the \emph{Shared Latent Space Hypothesis}: the subspaces \(\mathcal{S}^{(m)}\) are well-aligned across all \(m\), justifying a single \(U\).

---

\section{Shared Basis Construction Methods}

\subsection{Pretrained-Weights SVD}

Define
\[
\mathcal{X}_{\text{w}} = \left[ W_Q^{(1)} \mid W_K^{(1)} \mid W_V^{(1)} \mid \cdots \mid W_V^{(L)} \right].
\]
Calculate
\[
\mathcal{X}_{\text{w}} = U_{\text{w}} \Sigma_{\text{w}} V_{\text{w}}^\top,
\]
and take \(U = U_{\text{w}}^{(:,1:r)}\).

Then set \(A_*^{(\ell)} = U^\top W_*^{(\ell)}\).

\subsection{Covariance-Based SVD}

Estimate covariance
\[
C = \frac{1}{N} \sum_{i, \ell} (x_i^{(\ell)} - \mu)(x_i^{(\ell)} - \mu)^\top,
\]
of hidden states, with eigendecomposition
\[
C = U_{\text{c}} \Lambda U_{\text{c}}^\top.
\]
Take \(U = U_{\text{c}}^{(:,1:r)}\).

---

\section{Optimization Theory}

\subsection{Gradient Flow through Fixed Basis}

With
\[
W = U A,
\quad U^\top U = I_r,
\]
gradients w.r.t. \(A\) are
\[
\frac{\partial L}{\partial A} = U^\top \frac{\partial L}{\partial W}.
\]
Thus optimization acts as projected gradient descent onto the range of \(U\).

\subsection{Regularization and Generalization}

The low-rank constraint reduces effective parameter count, induces implicit regularization, and improves optimization conditioning.

\subsection{Migration vs From-Scratch Training}

Migration uses pretrained decomposition to initialize \(A\), facilitating fast convergence. From scratch involves learning \(U, A\), often less stable and with higher compute.

---

\section{Rank Sweep and Hardware Complexity}

\subsection{Parameter Count Reduction}

Original attention\_params:
\[
3 L d^2,
\]
SOA attention\_params:
\[
d r + 3 L r d = d r (1 + 3L).
\]

Relative ratio \(\approx \frac{37 r}{36 d}\) for BERT-base (\(L=12, d=768\)), resulting in \(4\% - 17\%\) reduction for ranks examined.

\subsection{FLOPs Reduction}

Original FLOPs per layer per projection \(\approx 3 T d^2\), SOA FLOPs
\[
\approx 4 T d r,
\]
achieving up to \(\approx 22\%\) of original projection cost at largest tested rank.

\subsection{M1/ANE Mapping}

Shared \(U\) enables reusing the large \(d \times r\) matrix efficiently on ANE; \(A_*^{(\ell)}\) matrices are small and easier to quantize and manage.

---

\section{Empirical Methodology (Outline)}

\begin{itemize}
\item Construct \(U\) using both pretrained-weights and covariance-based methods for ranks \(\{32,48,64,96,128\}\).
\item Initialize \(A_*^{(\ell)}\) from regression or pretrained weights.
\item Fine-tune only \(A_*^{(\ell)}\) (optionally with other parameters).
\item Evaluate GLUE/SQuAD tasks, long-context stability on M1, energy, latency.
\item Measure subspace alignment using principal angles.
\end{itemize}

---

\section{Summary}

Project SOA leverages classical numerical linear algebra and transformer theory to formulate a practical, theoretically well-founded approach to low-rank, shared-basis attention in BERT-base. The method is optimized for Apple M1 hardware enabling efficient long context attention. The fixed shared basis \(U\) is justified through PCA and subspace alignment theory. Optimization occurs over low-dimensional factors \(A_*^{(\ell)}\) with implicit projection, providing regularization and feasibility.

---

\section*{Acknowledgements}

This dossier compiles theoretical insights synthesized to support Project SOA initiatives for M1-optimized transformers.

\end{document}
